\begin{abstract}
\bi{LLM data agents learn what business metrics mean by exploring a database. Sharing the learned metric definitions across agents makes answers more accurate and consistent, but it turns each definition into long-lived state over a changing database. A definition's SQL is correct only while data properties that the SQL does not state still hold: a filtered key identifies one business record, joins do not fan out, and dates assign facts to the intended period. Data-only updates break these properties without any schema change; the shared query keeps running and silently returns wrong aggregates. We present \system, a middleware that maintains the validity of shared metric definitions and enforces it where they are used. \system\ derives four kinds of executable conditions from each learned definition and shows them sufficient for every business event to be counted once in its period. It keys their verdicts by dependency versions and admits a query that declares a definition only if the conditions hold on the snapshot on which the query runs, under a reuse rule that makes verdicts computed elsewhere sound for that snapshot. It repairs a broken definition only with the unique structural fix that passes regression, and withdraws the definition otherwise. Replaying \RpLibs\ definition libraries learned by three LLMs under \RpChanges\ data changes, maintained definitions never serve a wrong answer under the modeled changes, whereas schema-only maintenance serves \RpSchemaWrongModeled. Under concurrent writers, checking before execution answers on condition-violating data in up to \SnPreUnannMax\% of uses; snapshot-bound execution never does, at the same reader latency. Maintenance follows the distinct conditions an update touches and costs \CbStagCondSix\% of definition-level revalidation with 19 shared definitions. End to end, sharing lifts agent accuracy from \ScenAccNoShare\% to \DsNoguardAll\% and maintenance lifts it to \ScenAccCond\%; under the modeled breaking changes, maintenance raises accuracy from \DsNoguardModeled\% to \DsCondModeled\%, where a matched trajectory-retrieval baseline reaches \DsTrajModeled\%.}
{大语言模型数据智能体通过探索数据库学到业务指标的含义。在智能体之间共享学到的指标定义，能让答案更准确、口径更一致，但也使每个定义成为变化数据库上的长期状态。定义的 SQL 只在一些 SQL 本身没有写明的数据性质成立时才正确：过滤后的键对应一条业务记录，连接不放大行数，日期把事实归入预期的期间。纯数据更新可以在不改任何表结构的情况下破坏这些性质；共享查询照常执行，却静默返回错误的汇总值。本文提出 \system，一个维护共享指标定义的有效性、并在使用处强制这种有效性的中间层。\system\ 从每个学到的定义推出四类可执行条件，并证明它们足以保证每个业务事件在其期间内恰好被计数一次。它按依赖版本索引条件结论，只有当条件在查询执行所在的快照上成立时才放行声明使用该定义的查询，并给出一条复用规则，使在别处算出的结论对该快照同样可靠。它只用通过回归测试的唯一结构性修复来修复失效定义，否则撤下该定义。把 3 个大模型学到的 \RpLibs\ 个定义库在 \RpChanges\ 种数据变化下回放：在已建模的变化下，维护后的定义从不给出错误答案，而只看结构的维护给出了 \RpSchemaWrongModeled\ 个错误答案。存在并发写入时，执行前检查会让最多 \SnPreUnannMax\% 的使用在违反条件的数据上作答；绑定快照的执行一次也没有，读者延迟不变。维护工作沿更新触及的不同条件展开，19 个定义共享条件时，耗时为定义级重验证的 \CbStagCondSix\%。端到端实验中，共享把智能体正确率从 \ScenAccNoShare\% 提高到 \DsNoguardAll\%，维护进一步提高到 \ScenAccCond\%；在已建模的破坏性变化下，维护把正确率从 \DsNoguardModeled\% 提高到 \DsCondModeled\%，而匹配的轨迹检索基线为 \DsTrajModeled\%。}
\end{abstract}
\keywords{\bt{data agents, metric definitions, shared memory, validity conditions, snapshot isolation, database updates}{数据智能体，指标定义，共享记忆，有效性条件，快照隔离，数据库更新}}
\maketitle
